% GENERATED FILE. Edit canonical lessons, then run npm run generate:books.
% canonical-source-hash: fnv1a64:85a1d811c8f5c294
% canonical-lessons: RU-C499-naveshchat, RU-C499-gostepriimnyy, RU-C499-gostepriimstvo, RU-C499-chayepitiye, RU-C499-gostit

\chapter{Visiting}
\label{ch:ru-visiting}
\hlchaptermodality{499}

\begin{chapteropening}
\textbf{By the end of this chapter:} \emph{I can talk about visiting people and hospitality.}

Complete the last lesson of chapter 499: I can talk about visiting people and hospitality.
\end{chapteropening}

\section[\texorpdfstring{\ru{навещать} (naveshchát)}{навещать (naveshchát)}]{\ru{навещать} (naveshchát) — to visit (a person)}

\label{lesson:RU-C499-naveshchat}



\begin{quote}
Before the new one: say the Russian for someone of the same age, then the Russian for a classmate.
\end{quote}

\subsection*{You'll want to know: \ru{навещать}}
\textbf{\ru{навещать}} (\emph{naveshchát}) — \textquotedblleft{}to visit (a person)\textquotedblright{}.

\begin{quote}
\ru{По} \ru{воскресеньям} \ru{мы} \ru{навещаем} \ru{бабушку}.
\end{quote}

\begin{quote}
\emph{Po voskresényam my naveshcháyem bábushku.}
\end{quote}

\begin{quote}
On Sundays we visit Grandma.
\end{quote}

Imperfective; perfective \ru{навестить} (\ru{навещу}). Takes the accusative of the person.

\begin{cousinweb}
\emph{\ru{навещать}} is built on \emph{\ru{на}-} and the root of \emph{\ru{весть}} 'news'; EtymDB has no entry, and the link is usually made through \emph{\ru{навестить}}.
\end{cousinweb}

\begin{grammarlens}[title={What you've built}]
One more word for telling, explaining and describing at B1.
\end{grammarlens}

\subsection*{Your turn}
\begin{itemize}
\raggedright
  \item \emph{Say it:} \emph{naveshchát}
  \item \emph{Say it:} \emph{naveshchát}, once more
  \item \emph{Say it:} a sentence of your own with \emph{naveshchát}
  \item \emph{Recall:} say the Russian for someone of the same age, then the Russian for a classmate, then say \emph{naveshchát} again
\end{itemize}

\begin{tcolorbox}[breakable,colback=teal!4,colframe=teal!35!black,title={Before you move on}]
What does \ru{навещать} mean? (\textquotedblleft{}To visit (a person)\textquotedblright{}.) Say it once more.
\end{tcolorbox}
\section[\texorpdfstring{\ru{гостеприимный} (gostepriímnyy)}{гостеприимный (gostepriímnyy)}]{\ru{гостеприимный} (gostepriímnyy) — hospitable}

\label{lesson:RU-C499-gostepriimnyy}



\begin{quote}
Before the new one: say the Russian for a classmate, then the Russian for to visit.
\end{quote}

\subsection*{You'll want to know: \ru{гостеприимный}}
\textbf{\ru{гостеприимный}} (\emph{gostepriímnyy}) — \textquotedblleft{}hospitable\textquotedblright{}.

\begin{quote}
\ru{Хозяева} \ru{были} \ru{очень} \ru{гостеприимные}.
\end{quote}

\begin{quote}
\emph{Khozyáyeva býli óchen gostepriímnyye.}
\end{quote}

\begin{quote}
Our hosts were very hospitable.
\end{quote}

Short forms: \ru{гостеприимен}, \ru{гостеприимна}.

\begin{cousinweb}
\emph{\ru{гостеприимный}} is a compound of \emph{\ru{гость}} 'guest' and the old verb \emph{\ru{приимать}} 'to receive'; EtymDB has no entry.
\end{cousinweb}

\begin{grammarlens}[title={What you've built}]
One more word for telling, explaining and describing at B1.
\end{grammarlens}

\subsection*{Your turn}
\begin{itemize}
\raggedright
  \item \emph{Say it:} \emph{gostepriímnyy}
  \item \emph{Say it:} \emph{gostepriímnyy}, once more
  \item \emph{Say it:} a sentence of your own with \emph{gostepriímnyy}
  \item \emph{Recall:} say the Russian for a classmate, then the Russian for to visit, then say \emph{gostepriímnyy} again
\end{itemize}

\begin{tcolorbox}[breakable,colback=teal!4,colframe=teal!35!black,title={Before you move on}]
What does \ru{гостеприимный} mean? (\textquotedblleft{}Hospitable\textquotedblright{}.) Say it once more.
\end{tcolorbox}
\section[\texorpdfstring{\ru{гостеприимство} (gostepriímstvo)}{гостеприимство (gostepriímstvo)}]{\ru{гостеприимство} (gostepriímstvo) — hospitality}

\label{lesson:RU-C499-gostepriimstvo}



\begin{quote}
Before the new one: say the Russian for to visit, then the Russian for hospitable.
\end{quote}

\subsection*{You'll want to know: \ru{гостеприимство}}
\textbf{\ru{гостеприимство}} (\emph{gostepriímstvo}) — \textquotedblleft{}hospitality\textquotedblright{}.

\begin{quote}
\ru{Спасибо} \ru{за} \ru{ваше} \ru{гостеприимство}!
\end{quote}

\begin{quote}
\emph{Spasíbo za váshe gostepriímstvo!}
\end{quote}

\begin{quote}
Thank you for your hospitality!
\end{quote}

Neuter.

\begin{cousinweb}
\emph{\ru{гостеприимство}} is built on \emph{\ru{гостеприимный}} 'hospitable', a compound of \emph{\ru{гость}} 'guest' and \emph{\ru{приимать}} 'to receive'; EtymDB has no entry.
\end{cousinweb}

\begin{grammarlens}[title={What you've built}]
One more word for telling, explaining and describing at B1.
\end{grammarlens}

\subsection*{Your turn}
\begin{itemize}
\raggedright
  \item \emph{Say it:} \emph{gostepriímstvo}
  \item \emph{Say it:} \emph{gostepriímstvo}, once more
  \item \emph{Say it:} a sentence of your own with \emph{gostepriímstvo}
  \item \emph{Recall:} say the Russian for to visit, then the Russian for hospitable, then say \emph{gostepriímstvo} again
\end{itemize}

\begin{tcolorbox}[breakable,colback=teal!4,colframe=teal!35!black,title={Before you move on}]
What does \ru{гостеприимство} mean? (\textquotedblleft{}Hospitality\textquotedblright{}.) Say it once more.
\end{tcolorbox}
\section[\texorpdfstring{\ru{чаепитие} (chayepítiye)}{чаепитие (chayepítiye)}]{\ru{чаепитие} (chayepítiye) — a tea party, tea-drinking}

\label{lesson:RU-C499-chayepitiye}



\begin{quote}
Before the new one: say the Russian for hospitable, then the Russian for hospitality.
\end{quote}

\subsection*{You'll want to know: \ru{чаепитие}}
\textbf{\ru{чаепитие}} (\emph{chayepítiye}) — \textquotedblleft{}a tea party, tea-drinking\textquotedblright{}.

\begin{quote}
\ru{После} \ru{обеда} \ru{у} \ru{нас} \ru{всегда} \ru{чаепитие}.
\end{quote}

\begin{quote}
\emph{Pósle obéda u nas vsegdá chayepítiye.}
\end{quote}

\begin{quote}
After lunch we always sit down to tea.
\end{quote}

Neuter.

\begin{cousinweb}
\emph{\ru{чаепитие}} is a compound of \emph{\ru{чай}} 'tea' and \emph{\ru{питие}} 'drinking', from \emph{\ru{пить}} 'to drink'; EtymDB has no entry.
\end{cousinweb}

\begin{grammarlens}[title={What you've built}]
One more word for telling, explaining and describing at B1.
\end{grammarlens}

\subsection*{Your turn}
\begin{itemize}
\raggedright
  \item \emph{Say it:} \emph{chayepítiye}
  \item \emph{Say it:} \emph{chayepítiye}, once more
  \item \emph{Say it:} a sentence of your own with \emph{chayepítiye}
  \item \emph{Recall:} say the Russian for hospitable, then the Russian for hospitality, then say \emph{chayepítiye} again
\end{itemize}

\begin{tcolorbox}[breakable,colback=teal!4,colframe=teal!35!black,title={Before you move on}]
What does \ru{чаепитие} mean? (\textquotedblleft{}A tea party, tea-drinking\textquotedblright{}.) Say it once more.
\end{tcolorbox}
\section[\texorpdfstring{\ru{гостить} (gostít)}{гостить (gostít)}]{\ru{гостить} (gostít) — to stay as a guest}

\label{lesson:RU-C499-gostit}



\begin{quote}
Before the new one: say the Russian for hospitality, then the Russian for a tea party.
\end{quote}

\subsection*{You'll want to know: \ru{гостить}}
\textbf{\ru{гостить}} (\emph{gostít}) — \textquotedblleft{}to stay as a guest\textquotedblright{}.

\begin{quote}
\ru{Летом} \ru{я} \ru{гостила} \ru{у} \ru{тёти} \ru{в} \ru{деревне}.
\end{quote}

\begin{quote}
\emph{Létom ya gostíla u tyóti v derévne.}
\end{quote}

\begin{quote}
In the summer I stayed with my aunt in the village.
\end{quote}

Imperfective; present \ru{гощу}, \ru{гостишь} (\ru{ст} becomes \ru{щ}). \emph{\ru{Гостить} \ru{у} \ru{кого}}.

\begin{cousinweb}
\emph{\ru{гостить}} is built on \emph{\ru{гость}} 'guest'; EtymDB lists it, glossed 'to be on a visit, to stay', without a source.
\end{cousinweb}

\begin{grammarlens}[title={What you've built}]
One more word for telling, explaining and describing at B1.
\end{grammarlens}

\subsection*{Your turn}
\begin{itemize}
\raggedright
  \item \emph{Say it:} \emph{gostít}
  \item \emph{Say it:} \emph{gostít}, once more
  \item \emph{Say it:} a sentence of your own with \emph{gostít}
  \item \emph{Recall:} say the Russian for hospitality, then the Russian for a tea party, then say \emph{gostít} again
\end{itemize}

\begin{tcolorbox}[breakable,colback=teal!4,colframe=teal!35!black,title={Before you move on}]
What does \ru{гостить} mean? (\textquotedblleft{}To stay as a guest\textquotedblright{}.) Say it once more.
\end{tcolorbox}
